%%%%%%%%%%%%%%%%%%%%%%%%
\chapter{Design and Approach}
%%%%%%%%%%%%%%%%%%%%%%%%

The checker is implemented entirely within the Scala 3 compiler as a
\texttt{MiniPhase}, operating directly on the typed abstract syntax tree (AST)
produced by earlier compiler phases. It does not rely on any external solver,
rewriting engine, or verification framework. All information needed is obtained
through the compiler's own symbol and type APIs. This makes the checker
lightweight, and immediately available to any project that uses the Scala 3
compiler, without additional installation or configuration.

Furthermore, unlike proof assistants that enforce termination globally, this
checker is entirely opt-in. This aligns with the needs of a general-purpose
language like Scala, where the overhead of proving the termination of every
method is impractical.
The phase is enabled by the \texttt{-Ycheck-termination} compiler flag and only
methods annotated with \texttt{@terminates} and methods called by the annotated
methods are checked.

\section{Conservative, Sound Analysis}
The checker prioritizes soundness over completeness: it will never accept
non-terminating methods, however, it may reject methods that do in fact
terminate.
During the AST traversal, the checker assigns every value a \texttt{Size}
(either \texttt{Smaller}, \texttt{Same}, or \texttt{Unknown}) relative to the
parameters of the method, then, at each recursive call, the arguments to the
call are compared lexicographically to the method's parameters. However, the
sizes are inferred syntactically, meaning that the checker will reject programs
that terminate if the termination relies on semantic properties.

\section{Escape Hatches}
Because the analysis is conservative, rather than forcing developers to
restructure their code for the benefit of the checker, escape hatches are
provided as annotations, shifting the analysis on the developer. Misuse of these
escape hatches can cause the checker to accept a non-terminating function, so
they should be used with care.

The \texttt{@assumeTerminates} annotation applied to a method causes the checker
to accept any call to that method without inspecting its body. This is
especially useful for Scala's standard library methods whose bodies are
unavailable.

The \texttt{uncheckedTermination} method allows the user to wrap a block that is
not checked, this is useful for wrapping calls to external libraries when the
source is unavailable to the user.

The \texttt{@assumeDecreases} annotation applied to a method (typically a field
accessor or a selector-like method), tells the checker that the called method on
a value outputs a result that is strictly smaller than that value.
This is useful when it's semantically clear that the value decreases, but the
simple structural analysis fails.

The \texttt{@assumeWellFounded} annotation applied to a type bypasses the
well-foundedness check, asserting that the type admits no infinite descending
chains. It is needed for types that are structurally well-founded but whose
shape does not match the simple syntactic criterion the checker applies.

The \texttt{@decreases} annotation is not an escape hatch, but is used to
specifies in which order, should be used as the lexicographic measure for a
given method. This only affects the order in which arguments are compared, and
is necessary when the decreasing argument is not the first parameter, and the
default left-to-right order would cause the check to fail.


%%%%%%%%%%%%%%%%%%%%%%%%
\chapter{Implementation}
%%%%%%%%%%%%%%%%%%%%%%%%

The checker is implemented as a \texttt{MiniPhase} inside the
\texttt{dotty.tools.dotc.transform} package.
When a method is annotated with \texttt{@terminates}, is not abstract
(\texttt{Deferred}), and has not already been verified in a prior traversal, the
methods body is traversed. Successfully checked methods are accumulated in a
\texttt{checked} set shared across all invocations of the phase. A method that
has already been certified is not re-examined when it is called by another
annotated method. The phase is entirely inactive unless the
\texttt{-Ycheck-termination} compiler flag is set, ensuring zero overhead for
projects that do not opt in.

\section{Structural Decrease Analysis}
The core analysis tracks how the sizes of values evolve as execution proceeds.
A \texttt{valMap} associates local value symbols to a \emph{(source, size)}
pair, recording where the value came from and whether it is \texttt{Smaller}
or \texttt{Same} relative to that source. An analogous \texttt{varMap} tracks
mutable variables (\texttt{var}), updated at every assignment.

\subsection{Pattern Matching and ADTs}
Pattern matching on an Algebraic Data Type is the primary mechanism by which
structural decrease is established. When the selector of a \texttt{match} has
a synthetic \texttt{unapply} method, i.e., one generated automatically by the
compiler for a \texttt{case class} or \texttt{enum}, the checker knows that
the binders introduced by each pattern are strict sub-value of the
selector. The method \texttt{mapMatchedSymbols} walks the pattern tree. The
outermost binder (the variable that names the whole matched value) is recorded
as \texttt{Same} size as the selector, while any nested binders (the fields
extracted from the constructor) are recorded as \texttt{Smaller}.
If the \texttt{unapply} method is user-defined (not synthetic), the checker
cannot guarantee that deconstruction actually yields a smaller value, and
emits a warning, and does not consider the sub-value as \texttt{Smaller}.

\subsection{Field Access and \texttt{peelSelects}}
Beyond pattern matching, a decrease can also be established by selecting a
field of a value. The helper \texttt{peelSelects} resolves a chain of
field-access expressions back to an original symbol and determines whether a
decrease occurred along the way. A field selection is considered to decrease
the size of the value when the selected field's name corresponds to a
constructor parameter of a \texttt{case class}, or when the field's method is
annotated with \texttt{@assumeDecreases}.

\subsection{\texttt{methodOverrides} {\color{magenta}Find a better name}}
When a method is called on a value whose static type is a sealed trait or
abstract class, the checker cannot know at compile time which concrete
implementation will be invoked at runtime. To remain sound,
\texttt{methodOverrides} enumerates all sealed strict descendants of the type
and resolves the method name in each descendant's type. If any overrides are
found, the checker verifies each one independently. A call is accepted only if
every possible override terminates. If the type is concrete, no override
resolution is needed and the method symbol is used directly.

\subsection{Well-Foundedness}
A structural decrease is only a valid termination argument if the type being
decreased is well-founded, that is, it admits no infinite descending chain of
sub-values. The checker enforces this via \texttt{typeWellFounded}. A sealed
type is well-founded when all of its children are well-founded. A case object
(a leaf of an ADT) is trivially well-founded. A case class is well-founded
when all of its fields are stable members that correspond directly to
constructor parameters. If the type is not well-founded, the checker emits a
warning and considers the value of size \texttt{Unknown}.

\subsection{Lexicographic Ordering}
When checking a recursive call, the method \texttt{areSmaller} compares the list
of arguments against the list of parameters of the recursive method in order. It
calls \texttt{isLexicoDecreasing}, which steps through the argument–parameter
pairs. If the first pair is \texttt{Smaller}, the check succeeds provided the
type is well-founded. If it is \texttt{Same}, the check proceeds to the next
pair. If it is \texttt{Unknown}, the check fails. The ordering of the pairs can
be controlled by the \texttt{@decreases} annotation, which specifies an ordering
of parameters to use as the measure, overriding the default left-to-right order.


\section{Mutual Recursion}
The checker handles mutual recursion by inlineing called methods. A
\texttt{callStack} tracks the methods currently being analyzed. When a call to a
method is encountered, the checker first looks it up in \texttt{callStack}.
If it is already present, the call is recursive (possibly mutual) and
\texttt{areSmaller} is invoked to verify the decrease. If the method is not yet
on the stack, \texttt{traverseCalled} is invoked, which pushes the callee onto
the stack, maps the actual arguments to the callee's parameters in
\texttt{valMap}, traverses the callee's body, and then restores the stack and
map. This allows the checker to follow chains of mutual recursion of arbitrary
depth.


\section{While Loops}
While loops are treated separately from recursive calls. Since they rely on
mutable variable. The checker performs two traversals. In the first traversal,
call checking is disabled and \textttt{varMap} is reset to empty. The purpose of
this traversal is to determine whether there exists a variable that is assigned
a strictly smaller value (with respect to itself) inside the loop body. If such
a variable exists the loop terminates. Note that since well-founded types cannot
decrease indefinitely, we have two possible cases. Either the loop condition
becomes false, hence, the loop terminates. Or we hit a leaf of the ADT and the
assignment to a strictly smaller value results in an exception being thrown
(this is the case for e.g. \texttt{Nil.tail}) or the \texttt{@assumeDecreases}
escape hatch was wrongly used.
The checker then proceeds to a second traversal with call checking re-enabled,
and the original \texttt{varMap} restored so that any recursive calls inside the
loop body are themselves verified to be terminating.

\subsection{\texttt{varMap} Updates}
An \texttt{Assign} node records \texttt{Smaller} in \texttt{varMap} if
\texttt{areSmaller} can confirm the right-hand side is smaller than the current
value of the variable, and \texttt{Unknown} otherwise.

For branching and looping constructs, the checker explores each branch
independently and then merges the results where the worst-case size is kept.
Concretely, when traversing an \texttt{If} expression or \texttt{Match} cases,
we compute one \texttt{varMap} for each branch, then, for each variable in the
union of all maps, it is considered \texttt{Unknown} if \texttt{Unknown} in
any branch, \texttt{Smaller} if \texttt{Smaller} in all branches and
\texttt{Same} otherwise. Then the original \texttt{varMap} is updated with the
new information from the branching construct.
For \texttt{WhileDo}, the body might never execute, therefore, we merely merge
the original \texttt{varMap} with a \texttt{varMap} computed by traversing the
body of the loop.


%%%%%%%%%%%%%%%%%%%%
\chapter{Evaluation and Limitations}
%%%%%%%%%%%%%%%%%%%%
\section{\color{magenta} Performance and metrics}
\begin{itemize}
    \item Examples of when it works, and how.
    \item Examples of when it doesn't work, and why.
\end{itemize}


%%%%%%%%%%%%%%%%%%%%
\chapter{Correctness}
%%%%%%%%%%%%%%%%%%%%
\begin{itemize}
    \item Short correctness sketch, for recursion, mutual recursion, while
    loops, and, higher-order-function.
\end{itemize}


%%%%%%%%%%%%%%%%%%%%
\chapter{Discussion}
%%%%%%%%%%%%%%%%%%%%

